% !TeX root = ../sections/02_exercises.tex
\begin{exercise}
	数列
	\[
		a_n=(-1)^n+\frac{1}{n}
		\qquad
		(n=1,2,3,\ldots)
	\]
	は発散することを証明せよ。
\end{exercise}
\begin{background}
	この問題では，数列が収束しないことを示す。
	数列の発散を示す方法として，異なる極限をもつ部分列を見つける方法が有効である。

	\subsection{偶数番目と奇数番目に分けて調べる準備}

	\begin{theorem}[収束列の部分列]
		数列 $\{a_n\}$ が実数 $\alpha$ に収束するならば，
		その任意の部分列も $\alpha$ に収束する。
	\end{theorem}

	\begin{remark}
		今回の数列は
		\[
			a_n=(-1)^n+\frac{1}{n}
		\]
		である。

		偶数番目では
		\[
			a_{2m}=1+\frac{1}{2m}
		\]
		となり，これは $1$ に収束する。

		奇数番目では
		\[
			a_{2m-1}=-1+\frac{1}{2m-1}
		\]
		となり，これは $-1$ に収束する。
	\end{remark}
\end{background}
\begin{strategy}
	数列
	\[
		a_n=(-1)^n+\frac{1}{n}
	\]
	を偶数番目と奇数番目に分けて考える。

	偶数番目では
	\[
		a_{2m}=(-1)^{2m}+\frac{1}{2m}
		=
		1+\frac{1}{2m}
	\]
	であるから，
	\[
		a_{2m}\to 1
	\]
	である。

	一方，奇数番目では
	\[
		a_{2m-1}=(-1)^{2m-1}+\frac{1}{2m-1}
		=
		-1+\frac{1}{2m-1}
	\]
	であるから，
	\[
		a_{2m-1}\to -1
	\]
	である。

	もし元の数列 $\{a_n\}$ が収束するならば，
	すべての部分列は同じ極限に収束しなければならない。

	しかし，偶数番目の部分列は $1$ に収束し，
	奇数番目の部分列は $-1$ に収束する。
	これらは異なる極限である。

	したがって，元の数列 $\{a_n\}$ は収束しない。
	すなわち発散する。
\end{strategy}